% !PRINTED 475
\weberpage{475}
%<<CONT
närer linearer Substitutionen für den Modul 2 darstellen, die wir im \S.\,79 untersucht haben\footnote{Nach einer mündlichen Mittheilung ist dies der Weg, auf dem sie \gesperrt{Kronecker} gebildet hat.}.

Wir haben zu diesem Zweck die sieben Grössen durch drei Indices $(x_1, x_2, x_3)$ zu bezeichnen, die nach dem Modul 2 genommen sind und wobei die Combination $(0, 0, 0)$ ausgeschlossen ist. Man erhält so die sieben Grössen

\weq{1}{(1,0,0),\ (0,1,0),\ (0,0,1);\ (0,1,1),\ (1,0,1),\ (1,1,0),\ (1,1,1),}

und wenn man die $x_1, x_2, x_3$ durch die linearen Verbindungen

\weq{2}{a x_1 + b x_2 + c x_3,\ a_1 x_1 + b_1 x_2 + c_1 x_3,\ a_2 x_1 + b_2 x_2 + c_2 x_3}

ersetzt, worin die Substitution

\weq{3}{\begin{pmatrix} a, & b, & c \\ a_1, & b_1, & c_1 \\ a_2, & b_2, & c_2 \end{pmatrix}}

nach dem Modul 2 zu nehmen ist, so erhält man die Permutationsgruppe $P_{168}$.

Wenn wir, wie im \S.\,79, für $\tau$ die Substitution der Congruenzgruppe

\weq{4}{\tau = \begin{pmatrix} 1, & 0, & 1 \\ 1, & 0, & 0 \\ 0, & 1, & 0 \end{pmatrix}}

wählen, so können wir die Grössen $z$ so bezeichnen, dass sie durch Anwendung von $\tau$ und seinen Wiederholungen cyklisch in einander übergehen, wobei wir eine beliebige der Grössen (1) für $z_0$ wählen können, etwa so:

\weq{5}{z_0 = (1,0,0),\ z_1 = (1,1,0),\ z_2 = (1,1,1),\ z_3 = (0,1,1),\ z_4 = (1,0,1),\ z_5 = (0,1,0),\ z_6 = (0,0,1).}

Es ist dann $\omega$ so zu wählen, dass $z_0, z_4, z_6$ dadurch ungeändert bleiben und $z_1$ mit $z_2$, $z_3$ mit $z_5$ vertauscht werden. Dies giebt

\weq{6}{\omega = \begin{pmatrix} 1, & 0, & 0 \\ 0, & 1, & 0 \\ 0, & 1, & 1 \end{pmatrix},}

und daraus erhält man nach \S.\,72, (17):

\weq{7}{\chi = \begin{pmatrix} 1, & 0, & 0 \\ 0, & 0, & 1 \\ 0, & 1, & 1 \end{pmatrix}, \quad \Theta = \begin{pmatrix} 1, & 1, & 1 \\ 0, & 1, & 0 \\ 0, & 1, & 1 \end{pmatrix}.}
